\scanpage{69}
% archive sheet number: 68
\[
f(x,y,z)=x^2+y^2+z^2-\underbrace{(2+2(\alpha+\beta)+\alpha\beta)}_{\gamma=(\alpha+2)(\beta+2)-2}yz-\beta zx-\beta xy-(x+y+z)
\]
Les fonctions sous-quadratiques inv. par $\unsure{\Gamma|\mathscr{G}}$ sont donc les combinaisons linéaires de $f$ et de la constante $1$. Celles qui s'annulent sur les sommets du polyèdre, i.e. en $S_0=(0,0,0)$, sont les multiples de $f$.
\[
\left\{\begin{aligned}
\text{Matrice de }f_0\ \text{(\unsure{partie hom|partie homog|partie hom.})}&=\begin{pmatrix}
2&-\beta&-\beta\\-\beta&2&-\gamma\\-\beta&-\gamma&2
\end{pmatrix}\\
\delta'\ (\text{discr. divisé})&=4-\beta^2(\gamma+2)-\gamma^2\\
&=-(\alpha+2)(\beta+2)^2(\alpha+\beta)
\end{aligned}\right.
\]
\margin{Dans cas ico. $\alpha^2\unsure{+|-}\alpha-1=0$, $\beta=-1$ $\gamma=\alpha$, \uuline{$\delta'=1$} OK !}
\[
\left[\begin{aligned}
\text{inv. ssi }&\alpha\text{ partout }\ne\deleted{\unsure{+}}-2\\
&\beta\text{ partout }\ne-2\\
&\alpha+\beta:\text{ partout }\ne0
\end{aligned}\right]
\]
\margin{$\alpha=\unsure{2\cos\frac{\pi}{p},q|2\cos\frac{\pi}{p,q}}$}
où
\[
\gamma=2(1+\alpha+\beta)+\alpha\beta
\]
% margin note written sideways
\margin{\uline{NB} $f_0$ \unsure{quadratique} $\unsure{f}$ \unsure{bilinéaire|bilinéarisée|bilinéarisation}}
\[
\text{Matrice de }f=\begin{pmatrix}
2&-\beta&-\beta&-1\\
-\beta&2&-\gamma&-1\\
-\beta&-\gamma&2&-1\\
-1&-1&-1&0
\end{pmatrix}
\]
\[
\delta(f)=(\gamma+2)(\gamma-4\beta-6)=(\alpha+2)(\alpha-2)(\beta+2)^2
\]
\margin{$\delta(f)=\deleted{\unsure{4}}-\alpha-3$ $N_{\mathbb{Q}(\alpha)/\mathbb{Q}}\delta(f)=5$ !}
% Long downward arrow links the two notes in the right margin.
\diagram{\begin{tikzcd}[ampersand replacement=\&]
{}\arrow[d,bend left=15]\\{}
\end{tikzcd}}
il est \unsure{inv.} ssi $\alpha$ partout $\ne\pm2$ i.e. $\deleted{\unsure{4}}\unsure{2\cos\frac{\pi}{p}}$, $\beta$ partout $\ne-2$ $\unsure{2\cos\frac{\pi}{q}}$
\margin{(ici $p=\unsure{n}$ des sommets $q:$ \unsure{---} faces) (\uuline{\unsure{vrai}} $p,q$ \unsure{finis})}

\uline{L'unique} pt invariant par $\sigma_0,\sigma_1,\sigma_2$ i.e. par $\Gamma$ est
\[
o=\left(\frac{\deleted{\unsure{-}}\alpha}{2(\alpha+\beta)},\frac{-1}{2(\alpha+\beta)},\frac{-1}{2(\alpha+\beta)}\right)
\]
% slanted
\margin{\uline{si $2(\alpha+\beta)$ inv.}}
l'antipodisme est donné par
\[
\underline{a}(x,y,z)=\left(\frac{\alpha}{\alpha+\beta}-x,\frac{-1}{\alpha+\beta}-y,\frac{-1}{\alpha+\beta}-z\right)
\]
% slanted
\margin{si $\alpha+\beta$ inv.}
